\section{问题三：算力约束下的资源配置优化与结构性转移}

本章在总算力预算有限时，求使损失最小的参数量、数据量与质量配置，并分析预算跨量级时最优策略是否发生质变。目标函数取自问题二的广义标度律，约束为三部分算力之和不超预算。

\subsection{优化模型}
决策变量为参数量 $N>0$、数据量 $\Dt>0$ 与质量 $Q\in(0,1]$。$N,\Dt$ 分别以十亿参数、十亿 token 为单位；记 $U=10^9$，则物理计数为 $UN,U\Dt$。上下文长度 $\Lctx$ 为外生参数，配比在当前数值模型中固定。目标与 FLOPs 约束为
\begin{align}
\min_{N,\Dt,Q}\ & L(N,\Dt,Q)=E+A\,N^{-\alpha}+B\,\bigl(Q^{\theta}\Dt\bigr)^{-\beta},
\label{eq:q3-obj}\\
\text{s.t.}\ & \underbrace{6U^2N\Dt}_{\Ctrain}
+\underbrace{U\Dt\,[g(Q)-g(Q_0)]_+}_{C_Q}
+\underbrace{\eta U^2N\Dt\Lctx}_{\Cattn}\ \le\ C,
\label{eq:q3-budget}\\
& 0<Q\le 1,\qquad N>0,\ \Dt>0,
\label{eq:q3-dom}
\end{align}
其中基础训练开销取 Chinchilla 系数 $6$（每参数每 token 约 $6$ FLOPs，仅在标准 dense transformer 下成立；MoE 等稀疏架构需按激活比例折算，本文优化的模型族均按 dense 处理）。注意力开销是工程一阶近似，$\eta=2\times10^{-4}$；它刻画长上下文挤占预算的趋势，不作逐架构的精确 FLOPs 计数。$g(Q)$ 为每个物理 token 的质量成本（FLOPs/token），取指数、幂、对数三型之一；$[\cdot]_+$ 表示仅当质量超过基线 $Q_0$ 时才计增量成本。$Q_0=0.079$ 取自问题一域级质量中位数。显式写出 $U$ 可避免把“十亿”为单位的标度律变量直接代入物理 FLOPs 公式。

\subsection{KKT 解析刻画}
预算约束在最优处通常取等号，因为增大任一因素都降损失、必用尽预算。在 $Q>Q_0$ 的活跃段构造 Lagrange 函数
\begin{equation}
\mathcal L=E+A N^{-\alpha}+B(Q^{\theta}\Dt)^{-\beta}
+\mu\bigl(U^2N\Dt(6+\eta\Lctx)+U\Dt[g(Q)-g(Q_0)]-C\bigr),
\label{eq:q3-lagr}
\end{equation}
其一阶条件为
\begin{align}
\partial_N\mathcal L&=-\alpha A N^{-\alpha-1}+\mu U^2\Dt(6+\eta\Lctx)=0,
\label{eq:q3-kkt-n}\\
\partial_D\mathcal L&=-\beta B Q^{-\theta\beta}\Dt^{-\beta-1}
+\mu\bigl(U^2N(6+\eta\Lctx)+U[g(Q)-g(Q_0)]\bigr)=0,
\label{eq:q3-kkt-d}\\
\partial_Q\mathcal L&=-\theta\beta B Q^{-\theta\beta-1}\Dt^{-\beta}
+\mu U\Dt\,g'(Q)=0 .
\label{eq:q3-kkt-q}
\end{align}
由式~\eqref{eq:q3-kkt-n}--\eqref{eq:q3-kkt-d} 可消去乘子 $\mu$，得到参数量与数据量的边际平衡关系\cite{hoffmann2022}。质量条件须分段理解：$Q<Q_0$ 时提高质量不增加 $C_Q$ 却降低损失，故候选最优解不会停在该段；在 $Q=Q_0$ 处比较铰链两侧的单边导数，在 $Q_0<Q<1$ 的内点使用式~\eqref{eq:q3-kkt-q} 的边际收益等于边际成本条件，在 $Q=1$ 时检查上界条件。以上均为局部必要条件，不能单独证明全局最优。

\subsection{数值求解}
KKT 给出候选解的局部条件；数值上用序列二次规划配合多起点搜索\cite{sqp1995}。为避免变量尺度悬殊，对参数量与数据量取对数尺度 $(\ln N,\ln\Dt,Q)$，预算约束以非线性约束形式送入求解器；每档预算取不少于 $20$ 个初值，报告其中损失最低的可行解；多起点搜索不能证明全局最优。三档预算 $C\in\{10^{19},10^{22},10^{24}\}$，并在预算对数区间连续扫描以刻画转移；三型质量成本分别求解以比较其对最优解的影响。表~\ref{tab:q3_allocation} 汇总结果。

\textbf{可复现求解步骤。}
\begin{enumerate}[leftmargin=2em,itemsep=0pt,topsep=0.2em]
  \item 固定预算 $C$、成本函数 $g$ 与上下文长度；以式~\eqref{eq:q3-budget} 构造归一化可行余量 $(C-C_{\mathrm{tot}})/C\ge0$。
  \item 在 $(\ln N,\ln\Dt,Q)$ 空间取一个近似规模--数据平衡初值和其余随机初值，随机种子固定为 $42$；数值边界取 $N\in[10^{-3},10^5]$、$\Dt\in[10^{-2},10^6]$、$Q\in[10^{-3},1]$。三档主预算各用 $20$ 个初值，连续预算扫描用 $12$ 个。
  \item 对每个初值运行 SLSQP（最多 $500$ 次迭代，目标容差 $10^{-12}$）；只保留重算后满足 $C_{\mathrm{tot}}\le C(1+10^{-6})$ 的候选，按损失选最低者。该可行性筛选不等于全局最优证书。
\end{enumerate}
结构转移计算在 $\log_{10}C\in[18,25]$ 的 $36$ 个网格点上重复搜索，再对质量算力份额作数值差分；网格峰值应视作候选位置而非精确物理阈值。
以指数型质量成本和 $C=10^{22}$ FLOPs 为例，候选配置为约 $5.52$B 参数、$258.6$B token、$Q=0.968$。代回物理成本式，基础训练、质量与注意力开销分别约占预算的 $0.856$、$0.086$、$0.058$，三项之和约为 $1$。这一回代检验能发现单位错误或越预算解，但不能检验损失模型在该配置处的外推精度。

\begin{table}[H]
  \centering
  \caption{三档预算下的最优算力配置}
  \label{tab:q3_allocation}
  \small
  \resizebox{\textwidth}{!}{%
  \begin{tabular}{llrrrrrrr}
    \toprule
    $C$ (FLOPs) & $g(Q)$ & $N^*$ & $D^*$ & $Q^*$ & $L^*$ & $s_{train}$ & $s_Q$ & $s_{attn}$ \\
    \midrule
    $10^{19}$ & 指数型 & 0.193 & 7.9 & 0.5057 & 3.0729 & 0.9205 & 0.0167 & 0.0628 \\
     & 幂型 & 0.188 & 8.3 & 0.4879 & 3.0731 & 0.9361 & 0.0000 & 0.0639 \\
     & 对数型 & 0.188 & 8.3 & 0.4879 & 3.0731 & 0.9361 & 0.0000 & 0.0639 \\
    & Chinchilla 基线 & 1.249 & 1.2 & 0.4879 & 3.3066 & 0.9361 & 0.0000 & 0.0639 \\
    \midrule
    $10^{22}$ & 指数型 & 5.483 & 261.4 & 0.9665 & 2.1507 & 0.8600 & 0.0813 & 0.0587 \\
     & 幂型 & 5.758 & 240.3 & 1.0000 & 2.1525 & 0.8300 & 0.1133 & 0.0567 \\
     & 对数型 & 5.225 & 287.8 & 1.0000 & 2.1458 & 0.9023 & 0.0361 & 0.0616 \\
    & Chinchilla 基线 & 39.499 & 39.5 & 0.4879 & 2.2813 & 0.9361 & 0.0000 & 0.0639 \\
    \midrule
    $10^{24}$ & 指数型 & 40.740 & 3773.9 & 1.0000 & 1.9139 & 0.9225 & 0.0145 & 0.0630 \\
     & 幂型 & 40.893 & 3747.8 & 1.0000 & 1.9140 & 0.9195 & 0.0177 & 0.0628 \\
     & 对数型 & 40.280 & 3854.6 & 1.0000 & 1.9136 & 0.9316 & 0.0048 & 0.0636 \\
    & Chinchilla 基线 & 394.989 & 395.0 & 0.4879 & 1.9935 & 0.9361 & 0.0000 & 0.0639 \\
    \bottomrule
  \end{tabular}
  }
\end{table}


预算增大时最优损失单调下降（指数型下从 $3.098$ 经 $2.151$ 降至 $1.914$），逐步趋近不可约损失 $1.69$，参数量与数据量单调增，预算利用率约 $100\%$（约束活跃）；相对质量固定为基线、规模等分的 Chinchilla 基线方案，最优配置的损失改善约 $0.15$ 到 $0.58$（表~\ref{tab:q3_allocation}）。三型质量成本的最优解在中高预算趋于一致，仅在低预算档因质量成本经济性不同而分化：对数型使高质量相对可得，故低预算即把质量推到边界；指数型使高质量昂贵，故低预算下质量停在内点。三档预算的算力份额构成见图~\ref{fig:q3_allocation}。

\begin{figure}[H]
  \centering
  \includegraphics[width=0.9\textwidth,height=0.8\textheight,keepaspectratio]{fig_q3_allocation_stacked.pdf}
  \caption{三档预算下的算力份额构成}
  \label{fig:q3_allocation}
\end{figure}

基础训练份额始终主导（$0.80$--$0.92$），注意力份额随预算缓升，而质量提升份额在所示三档预算中随预算增大回落（图~\ref{fig:q3_allocation}）。这一份额变化提示资源结构可能转移，需要沿连续预算扫描定位。

\subsection{结构性转移的定义与识别}
给出可判定定义：以各部分算力份额 $s_X(C)=C_X^\ast/C$ 对预算对数的导数峰值定位转移点，
\begin{equation}
C^{\dagger}=\arg\max_{C}\ \left|\frac{\mathrm d\,s_X(C)}{\mathrm d\,\log C}\right| ,
\label{eq:q3-transition}
\end{equation}
并检查数值解是否触及质量边界或质量成本铰链。连续扫描预算对数区间、数值差分求份额导数可给出候选转移点；离散网格和成本函数选择会影响其位置。

\textbf{结果。}在当前 $\theta=0.5$、成本函数和扫描网格设定下，质量份额导数在预算约 $2.5\times10^{22}$ FLOPs 处出现峰值，作为模型内的候选结构转移点。基线质量 $0.079$ 较低，质量成本铰链在较低预算即已激活；这一结果同时依赖经验基线和假定的质量指数，不能称为纯数据驱动结论。转移点见图~\ref{fig:q3_transition}。

\begin{figure}[H]
  \centering
  \includegraphics[width=0.9\textwidth,height=0.8\textheight,keepaspectratio]{fig_q3_transition_line.pdf}
  \caption{最优份额演化与结构性转移点}
  \label{fig:q3_transition}
\end{figure}

份额随预算对数演化曲线上叠加的份额导数，其峰值位置即由式~\eqref{eq:q3-transition} 定位的转移点（图~\ref{fig:q3_transition}）。由于该最优轨迹来自确定性优化、不存在重复抽样，图中不画置信带，而以份额导数曲线的峰值给出结构判据，当前实现未计算 KKT 残差，不能用一阶条件数值精度为转移位置背书。

\subsection{上下文长度的临界值与敏感性}
上下文长度为外生参数。令注意力开销等于训练开销 $\eta N\Dt\Lctx=6N\Dt$，与参数量、数据量无关，解得临界上下文
\begin{equation}
\Lctx^{\mathrm{crit}}=\frac{6}{\eta}=\frac{6}{2\times10^{-4}}=3\times10^{4}\ \text{tokens}.
\label{eq:q3-lctxcrit}
\end{equation}
当上下文长度达到该值时注意力开销与训练开销相当，超过它长上下文将显著挤压训练预算。敏感性扫描的取值严格取自模型架构元数据的实测集，即 $2048$、$4096$、$8192$、$32768$ 与 $131072$。图~\ref{fig:q3_lctx} 展示最优配置随上下文长度的变化。

\begin{figure}[H]
  \centering
  \includegraphics[width=0.92\textwidth,height=0.8\textheight,keepaspectratio]{fig_q3_lctx_panels.pdf}
  \caption{最优配置随上下文长度的变化}
  \label{fig:q3_lctx}
\end{figure}

在中档预算下，上下文由 $2048$ 增至 $131072$ 时，注意力份额从约 $0.058$ 升至约 $0.776$（图~\ref{fig:q3_lctx}）。临界值 $3\times10^{4}$ 位于 $8192$ 与 $32768$ 之间；上下文不小于 $32768$ 时，按当前成本模型，注意力开销超过基础训练开销，压缩了可用于参数量和数据量的预算。可行域见图~\ref{fig:tikz_q3_feasible}。

\begin{figure}[H]
  \centering
  \includegraphics[width=0.8\textwidth,height=0.8\textheight,keepaspectratio]{tikz_q3_feasible.pdf}
  \caption{算力预算可行域与结构转移}
  \label{fig:tikz_q3_feasible}
\end{figure}

在参数量--数据量--质量空间中，由式~\eqref{eq:q3-budget} 界定的预算面、等成本线与最优点位置一同表明最优解落在预算面上（约束活跃），且随预算增大最优点沿预算面移动时穿过结构转移分区（图~\ref{fig:tikz_q3_feasible}）。综合本章：预算增大使损失单调改善并趋近不可约损失，最优策略在预算约 $2.5\times10^{22}$ 处发生质量投入结构的转移，长上下文以临界值 $3\times10^{4}$ 为界挤压训练预算。
